\documentclass{article}
\usepackage[utf8]{inputenc}
\usepackage{amsmath,amssymb,amsthm,mathtools}
\usepackage{algorithm,algorithmic}
\usepackage{graphicx}
\usepackage{hyperref}
\usepackage{microtype}
\usepackage{url}
\usepackage{booktabs}
\usepackage{natbib}
\usepackage[margin=1in]{geometry}
\usepackage{xcolor}

\DeclareMathOperator{\argmin}{arg\,min}
\DeclareMathOperator{\argmax}{arg\,max}

\newtheorem{theorem}{Theorem}
\newtheorem{proposition}{Proposition}
\newtheorem{lemma}{Lemma}
\newtheorem{remark}{Remark}
\newtheorem{corollary}{Corollary}

\title{Block-Wise Differentiable Sinkhorn Attention:\\Tail-Refinement Gradients with a Gap-Aware Dustbin Bridge}

\author{
Dylan B. Forde \\
\texttt{forde.dylan@gmail.com}
}

\date{}

\begin{document}

\maketitle

\begin{abstract}
We study long-context balanced entropic optimal transport (OT) attention on TPU hardware. The central object is a stopped-base, fixed-depth tail-refinement surrogate: after a stopped $T$-step Sinkhorn solve, we unroll a short refinement tail and differentiate that surrogate exactly. For the production kernel we use the $R=2$ case. Its backward pass forms a staircase of four transport plans, and we show that all non-base plans can be reconstructed exactly from a single base plan by row and column rescaling. This yields a block-wise implementation with fixed-bandwidth cost $O((T+R)LW)$, $O(Ld)$ input storage, and $O(L)$ additional HBM usage for fixed head dimension $d$ and band width $W$.

The first gap-aware extension is deliberately minimal and theorem-aligned: the current \texttt{dustbin\_block} path is formalized as the same balanced surrogate on an augmented state space, so the exact $R=2$ factorization lifts directly to the transport path exercised by the TPU runs. We also record a local surrogate-bias bound under fixed-support contraction assumptions. On synthetic masked problems, the optimized kernel matches exact autodiff of the same centered surrogate to within $10^{-5}$--$10^{-10}$ depending on tensor and length. On TPU v6e-8, the short Pfam screen completes end-to-end for the four core balanced and dustbin-augmented configurations, and a promoted balanced $R=2$ run sustains roughly $8.5$ examples per second through a three-hour budget, reaching step $1437$ with repeated checkpoint synchronization. Together these results support a single paper story: exact fixed-depth backward theory, a theorem-matching gap-aware bridge, and practical long-budget trainability of the production path.
\end{abstract}

\section{Introduction}

Pairwise alignment and structural matching remain central problems for sequence modeling. Classical dynamic-programming methods such as Smith--Waterman \citep{smith1981identification} provide optimal alignments but are not differentiable, so they cannot be inserted directly into end-to-end neural architectures. Entropic optimal transport (OT) offers a natural differentiable relaxation \citep{cuturi2013sinkhorn, peyre2019computational}: it replaces a discrete assignment with a smooth transport plan computed by Sinkhorn scaling.

For long-context models, however, two bottlenecks dominate. First, dense OT attention materializes $L\times L$ scores or plans, which is prohibitive at large $L$. Second, backward-pass design is difficult. Exact autodiff through many Sinkhorn steps is expensive and memory-hungry, while implicit differentiation can require large and unstable linear solves \citep{luise2018differential, bai2019deep}. The key question is therefore not only how to make OT attention differentiable, but how to make its differentiated object explicit, exact for the surrogate actually used, and compatible with accelerator kernels.

We address this with a stopped-base, fixed-depth \emph{tail-refinement surrogate}. After a stopped $T$-step balanced Sinkhorn solve, we append a short differentiable tail of depth $R$ and differentiate that surrogate exactly. The production TPU path uses $R=2$. The resulting backward pass forms a staircase of four transport plans, and the signature theorem of the paper is that all non-base plans in that staircase can be recovered exactly from one materialized base plan by row and column rescaling. This turns a seemingly multi-plan backward pass into a one-plan tile computation plus cheap vector modifiers, which maps well onto block-wise TPU kernels.

\paragraph{Gap-aware bridge.}
The central theorem is balanced and fixed-support. The first extension carried in the unified manuscript is intentionally narrow: the current \texttt{dustbin\_block} path is not treated as a new OT solver, but as the same balanced surrogate on an augmented state space. This keeps the gap-aware transport path inside the algebraic orbit of the main theorem, rather than forcing a new theorem prematurely.

\paragraph{Related context.}
Our work is closest in spirit to recent efforts to make OT-style attention practical at accelerator scale, including FlashSinkhorn \citep{flashsinkhorn2024}. It is also complementary to long-context transformer architectures that obtain efficiency through structured sparsity rather than transport constraints, such as Longformer \citep{beltagy2020longformer} and BigBird \citep{zaheer2020bigbird}. The point of departure here is that we focus on an exact fixed-depth backward factorization for Sinkhorn attention itself, rather than on sparse softmax attention or on an implicit OT solve.

The paper makes three main contributions:
\begin{itemize}
    \item an explicit balanced fixed-support tail-refinement surrogate together with its exact $R=2$ adjoint and single-base-plan factorization;
    \item a theorem-aligned dustbin augmentation showing that the current gap-aware transport path is the same balanced surrogate on an enlarged state space;
    \item an empirical bridge from exactness validation to TPU-scale training, comprising a completed four-configuration Pfam screen and a promoted long-budget balanced run.
\end{itemize}

\section{Method}

\subsection{Forward Pass: Block-Wise Balanced Sinkhorn}

Let $Q,K,V$ be query, key, and value tensors with head dimension $d$, and let $\Omega$ denote a fixed active support with band width $W$. We absorb the entropic temperature $\varepsilon$ into the score matrix
\[
S_{ij} = \frac{Q_i K_j^\top}{\sqrt{d}\,\varepsilon}.
\]
The stopped base solve and the refinement tail both use the same masked balanced half-steps
\begin{align}
    u_i^{(t+1)} &= -\log\!\Bigl[\sum_{j:(i,j)\in\Omega} \exp(S_{ij} + v_j^{(t)})\Bigr], \\
    v_j^{(t+1)} &= -\log\!\Bigl[\sum_{i:(i,j)\in\Omega} \exp(S_{ij} + u_i^{(t+1)})\Bigr].
\end{align}
The implementation computes these updates in tiles, in the same spirit as IO-aware attention kernels. For fixed band width $W$, a $T$-step base solve followed by an $R$-step tail has cost $O((T+R)LW)$ time, $O(Ld)$ input storage, and $O(L)$ additional HBM usage.

\begin{remark}[Fixed-Budget \texorpdfstring{$\varepsilon$}{epsilon}-Scaling]
The implementation optionally uses a short continuation schedule $\varepsilon_1 \ge \cdots \ge \varepsilon_S = \varepsilon$ during the stopped base solve. If the stage iteration counts satisfy $\sum_{s=1}^{S} T_s = T$, the asymptotic cost remains $O(TLW)$ time and $O(L)$ additional HBM memory for fixed head dimension.
\end{remark}

\subsection{Balanced Tail Refinement and the \texorpdfstring{$R=2$}{R=2} Staircase}

Let $(u^{(0)},v^{(0)})$ be the stopped potentials produced by the base $T$-step solve. The differentiated object is not the full $T$-step map; it is the fixed-depth surrogate
\[
\widetilde O^{[R]}(Q,K,V;u^{(0)},v^{(0)}) = \mathcal{A}(Q,K,V,u^{(R)},v^{(R)}),
\]
where the refinement iterates are recomputed from the stopped base pair. The implementation stores centered dual representatives for numerical stability, while the algebra below uses ungauged representatives; Appendix~\ref{app:proofs} proves that the centered and ungauged forward recurrences are gauge-equivalent.

Define staircase plans
\[
P_{ij}^{(a,b)} =
\begin{cases}
\exp(S_{ij}+u_i^{(a)}+v_j^{(b)}), & (i,j)\in\Omega, \\
0, & (i,j)\notin\Omega.
\end{cases}
\]
Let $\mathcal{L}$ be a scalar loss depending on $\widetilde O^{[R]}$, let $G=\partial \mathcal{L}/\partial \widetilde O^{[R]}$, and define
\[
Z_{ij}=\langle G_i, V_j\rangle,
\qquad
g_u=(P^{(R,R)}\odot Z)\mathbf{1},
\qquad
g_v=(P^{(R,R)}\odot Z)^\top\mathbf{1}.
\]

\begin{proposition}[Exact Fixed-Depth Tail Adjoint]\label{prop:tail_adjoint_main}
The cotangents of the refinement duals satisfy
\begin{align}
    \bar v^{(R)} &= g_v, \\
    \bar u^{(R)} &= g_u - P^{(R,R)}\bar v^{(R)}, \\
    \bar v^{(t-1)} &= -\bigl(P^{(t,t-1)}\bigr)^\top \bar u^{(t)}, \qquad t=R,R-1,\dots,1, \\
    \bar u^{(t-1)} &= -P^{(t-1,t-1)}\bar v^{(t-1)}, \qquad t=R,R-1,\dots,2,
\end{align}
where the final line stops at $t=2$ because the base state is stop-gradient. The score-space cotangent is
\begin{equation}
\label{eq:generic_score_vjp_unified}
\bar S
=
P^{(R,R)} \odot Z
-
\sum_{t=1}^R
\Bigl[
P^{(t,t)} \odot \bigl(\mathbf{1}\bar v^{(t)\top}\bigr)
+
P^{(t,t-1)} \odot \bigl(\bar u^{(t)}\mathbf{1}^\top\bigr)
\Bigr].
\end{equation}
\end{proposition}

For the production kernel we use $R=2$, which yields the staircase
\[
\bigl(P^{(2,2)},\,P^{(2,1)},\,P^{(1,1)},\,P^{(1,0)}\bigr).
\]

\begin{theorem}[Single-Base-Plan Reconstruction]\label{thm:single_base_main}
For any staircase pair $(a,b)\in\{(2,2),(2,1),(1,1),(1,0)\}$,
\[
P^{(a,b)}
=
P^{(2,2)}
\odot
\exp(u^{(a)}-u^{(2)})
\odot
\exp(v^{(b)}-v^{(2)}).
\]
Consequently, the exact $R=2$ score-space VJP can be evaluated while materializing only one transport tile per block pair.
\end{theorem}

This theorem is the conceptual center of the manuscript. It converts the explicit four-plan $R=2$ backward into a one-plan materialization plus cheap row and column modifiers, which is why the method maps cleanly to TPU kernels.

\subsection{Gap-Aware Dustbin Bridge}

The current gap-aware production path does not use a new OT solver. Instead, it augments the balanced problem. Let the base sequence length be $L$ and let $B=128$ denote the dustbin block size used in the implementation. Define augmented tensors
\[
\widetilde Q,\widetilde K,\widetilde V \in \mathbb{R}^{(L+B)\times d}
\]
by appending one learned dustbin token per side together with $B-1$ inactive filler tokens. The masks are augmented analogously: one appended token is active, the rest are inactive fillers. The implementation then applies the same balanced tail-refinement solver to this augmented system, with effective band width
\[
\widetilde W = \max(W,L).
\]

\begin{theorem}[Dustbin Augmentation Theorem]\label{thm:dustbin_aug}
Let $(\widetilde Q,\widetilde K,\widetilde V,\widetilde\Omega)$ denote the dustbin-augmented tensors and active support constructed above. Then the current \texttt{dustbin\_block} path is exactly the same balanced fixed-support tail-refinement surrogate as in Section~2.2, but applied on the augmented state space.
\end{theorem}

\begin{corollary}[Exact \texorpdfstring{$R=2$}{R=2} Adjoint on the Augmented State Space]\label{cor:dustbin_r2}
Apply Proposition~\ref{prop:tail_adjoint_main} and Theorem~\ref{thm:single_base_main} to the augmented system. Then the current dustbin-augmented transport path inherits the exact fixed-depth adjoint and the same single-base-plan reconstruction on the augmented support.
\end{corollary}

This is the first gap-aware extension carried into the unified paper because it already matches the cloud path. It is a bridge theorem, not a departure from the central balanced factorization.

\subsection{A Local Surrogate-Bias Bound}

The fixed-depth backward is exact for the surrogate it differentiates. A separate question is how far that surrogate gradient can be from the gradient through the full stopped-base computation. Let
\[
z=(u,v),\qquad \Phi_x(z)=z^+,
\]
be the centered one-step refinement map at differentiable parameters $x$, and let $A_x(z)$ denote the transport application map. Write
\[
z_T(x)=\Phi_x^T(z_0),
\qquad
\Psi_x^R(z)=A_x\!\bigl(\Phi_x^R(z)\bigr).
\]
Then the full fixed-depth gradient is
\[
g_{\mathrm{full}}^{T,R}(x)=\nabla_x \Psi_x^R(z_T(x)),
\]
while the stop-gradient surrogate uses
\[
g_{\mathrm{surr}}^{T,R}(x)=\partial_x \Psi_x^R(z)\big|_{z=z_T(x)}.
\]

\begin{theorem}[Local Surrogate-Bias Bound]\label{thm:surrogate_bias}
Assume fixed support, $C^1$ dependence of $\Phi_x$ and $A_x$ on $(x,z)$, and constants $\rho\in[0,1)$, $L_\Phi>0$, and $L_A>0$ such that
\[
\|D_z\Phi_x(z)\|\le\rho,
\qquad
\|D_x\Phi_x(z)\|\le L_\Phi,
\qquad
\|D_z A_x(z)\|\le L_A
\]
on a forward-invariant neighborhood. Then
\[
g_{\mathrm{full}}^{T,R}(x)-g_{\mathrm{surr}}^{T,R}(x)
=
D_z\Psi_x^R(z_T(x))\,D_x z_T(x),
\]
and
\[
\bigl\|g_{\mathrm{full}}^{T,R}(x)-g_{\mathrm{surr}}^{T,R}(x)\bigr\|
\le
\frac{L_A L_\Phi}{1-\rho}\,\rho^R.
\]
\end{theorem}

The point of this bound is modest but useful: once the support is fixed and the refinement map is locally contractive, the stop-gradient bias decays geometrically in the tail depth $R$. Appendix~\ref{app:bias} gives the short proof.

\section{Validation and Experiments}

\subsection{Exactness Validation of the Masked Surrogate}

We first separate surrogate correctness from systems scale. On synthetic masked problems for which exact autodiff through the same centered surrogate remains tractable, we compare the optimized $R=2$ kernel against a pure-JAX reference. Table~\ref{tab:exactness_unified} reports representative maximum absolute gradient deviations under the default validation setting ($\varepsilon=1$, band width $256$, $R=2$).

\begin{table}[h]
    \centering
    \begin{tabular}{@{}rrrr@{}}
        \toprule
        Length $L$ & $\max |\Delta Q|$ & $\max |\Delta K|$ & $\max |\Delta V|$ \\
        \midrule
        $512$ & $1.05 \times 10^{-5}$ & $1.73 \times 10^{-6}$ & $1.06 \times 10^{-9}$ \\
        $1024$ & $5.70 \times 10^{-6}$ & $8.13 \times 10^{-7}$ & $2.55 \times 10^{-10}$ \\
        $2048$ & $5.50 \times 10^{-6}$ & $9.17 \times 10^{-7}$ & $2.29 \times 10^{-10}$ \\
        \bottomrule
    \end{tabular}
    \caption{Maximum absolute gradient deviations between the optimized masked $R=2$ kernel and exact autodiff through the same centered surrogate.}
    \label{tab:exactness_unified}
\end{table}

Additional forward-consistency checks run through lengths $512$, $1024$, $2048$, $4096$, $8192$, and $16384$, with extended spot checks at $32768$, $65536$, and $131072$. We treat the larger checks as forward-consistency evidence rather than as exact-gradient checks.

\subsection{Short Pfam Screen on TPU v6e-8}

The first empirical bridge is the short Pfam screen on TPU v6e-8. All four core configurations complete end-to-end:
\[
\texttt{k2\_eps1},\quad
\texttt{k4\_eps1},\quad
\texttt{muot\_dustbin\_mono},\quad
\texttt{muot\_dustbin\_hybrid}.
\]
Each uses $50{,}000$ training pairs, $T=15$ base iterations, band width $1024$, and global batch size $64$. Table~\ref{tab:screen_unified} reports the step-50 screen snapshot.

\begin{table}[h]
    \centering
    \begin{tabular}{@{}lcccc@{}}
        \toprule
        Config & Step 50 Loss & Step 50 CE & Step 50 Mono & Throughput (ex/s) \\
        \midrule
        \texttt{k2\_eps1} & 3.5778 & 5.3223 & 0.3353 & 8.46 \\
        \texttt{k4\_eps1} & 3.5779 & 5.3223 & 0.3353 & 8.38 \\
        \texttt{muot\_dustbin\_mono} & 3.9030 & 5.3258 & 0.4155 & 8.41 \\
        \texttt{muot\_dustbin\_hybrid} & 4.6250 & 5.3244 & 0.4162 & 8.45 \\
        \bottomrule
    \end{tabular}
    \caption{Short TPU v6e-8 Pfam screen across the four core balanced and dustbin-augmented configurations.}
    \label{tab:screen_unified}
\end{table}

Two immediate conclusions follow. First, on the balanced path the short $R=2$ tail is empirically indistinguishable from $R=4$ at this screening budget, while remaining slightly faster. Second, the theorem-aligned dustbin bridge is operational at the same systems scale as the balanced path.

\subsection{Promoted Balanced Run}

We then promoted the balanced $\texttt{k2\_eps1}$ configuration to the longer-budget TPU training path. This promoted run keeps the screened objective weights,
\[
\lambda_{\mathrm{rec}}=1,\qquad \lambda_{\mathrm{ce}}=0,\qquad \lambda_{\mathrm{mono}}=0,
\]
so the reconstruction loss is optimized directly while CE and monotonicity are tracked diagnostically. The run sustained approximately $8.5$ examples per second through the three-hour budget, repeatedly synchronized checkpoints to GCS, and reached step $1437$. Table~\ref{tab:promoted_unified} shows representative snapshots from the run.

\begin{table}[h]
    \centering
    \begin{tabular}{@{}rrrr@{}}
        \toprule
        Step & Loss & CE & Throughput (ex/s) \\
        \midrule
        $0$ & 2.9320 & 5.4281 & 2.67 \\
        $500$ & 1.8221 & 5.2597 & 8.50 \\
        $1000$ & 0.9801 & 5.2754 & 8.52 \\
        $1400$ & 1.2839 & 5.2895 & 8.52 \\
        \bottomrule
    \end{tabular}
    \caption{Representative checkpoints from the promoted balanced TPU run.}
    \label{tab:promoted_unified}
\end{table}

This promoted run should be interpreted narrowly but positively. It is strong evidence that the theorem-backed balanced $R=2$ path is stable, efficient, and checkpointable under a realistic TPU budget. It is not by itself a claim that sparse alignment CE was the optimized quantity in this specific promotion, because CE was logged diagnostically rather than weighted into the training objective.

\section{Discussion}

The unified picture is now clear. The core contribution remains the balanced fixed-support tail-refinement surrogate together with the exact $R=2$ single-base-plan factorization. The dustbin bridge matters because it keeps the first gap-aware extension inside the same theorem family and already matches the current cloud path. The empirical evidence aligns with that theory stack: exactness holds on small masked problems, the four core balanced and dustbin-augmented screen configurations complete on TPU, and the promoted balanced path trains cleanly through a long budget.

What remains outside the present theorem spine should stay outside it. In particular, a genuinely KL-unbalanced transport theorem, differentiable support geometry, and multimodal routing all remain reasonable follow-ons, but they are not required for the central claim of this manuscript.

\section{Conclusion}

We presented a single-base-plan backward theory for long-context Sinkhorn attention and carried it through the first gap-aware extension that already matches the production code path. The balanced fixed-support $R=2$ factorization remains the centerpiece. The dustbin augmentation theorem lifts that factorization onto the current gap-aware transport path without changing the underlying surrogate. Exactness validation, a completed four-way TPU screen, and a promoted long-budget balanced run together show that the theorem is not merely algebraic: it maps to a practical TPU training path.

\bibliographystyle{plainnat}
\bibliography{references}

\newpage
\appendix

\section{Balanced Tail-Refinement Proofs}
\label{app:proofs}

\paragraph{Notation.}
For the appendix proofs we work with ungauged representatives satisfying the raw masked logsumexp half-step equations. Let $\Omega$ denote the active support and define
\[
P^{(a,b)}_{ij} =
\begin{cases}
\exp(S_{ij}+u_i^{(a)}+v_j^{(b)}), & (i,j)\in\Omega, \\
0, & (i,j)\notin\Omega.
\end{cases}
\]

\begin{align}
    u_i^{(t+1)} &= -\log\!\Bigl[\sum_{j:(i,j)\in\Omega} \exp(S_{ij} + v_j^{(t)})\Bigr], \label{eq:u_update_app_unified}\\
    v_j^{(t+1)} &= -\log\!\Bigl[\sum_{i:(i,j)\in\Omega} \exp(S_{ij} + u_i^{(t+1)})\Bigr]. \label{eq:v_update_app_unified}
\end{align}

\paragraph{Centering map.}
Let $\mathbf{1}_q$ and $\mathbf{1}_k$ denote the all-ones vectors on valid queries and keys, and let
\[
m = \frac{\mathbf{1}_q}{\mathbf{1}_q^\top\mathbf{1}_q}.
\]
The implementation stores centered representatives via
\[
\mathcal{C}(u,v)=\bigl(u-(m^\top u)\mathbf{1}_q,\; v+(m^\top u)\mathbf{1}_k\bigr).
\]

\begin{lemma}[Centering Pullback]
If $(\bar u,\bar v)$ are cotangents on the outputs of $\mathcal{C}$, then
\[
\mathcal{C}^*(\bar u,\bar v)
=
\Bigl(
\bar u + m\bigl(\mathbf{1}_k^\top \bar v - \mathbf{1}_q^\top \bar u\bigr),
\;
\bar v
\Bigr).
\]
\end{lemma}

\begin{proof}
Write $s=m^\top u$, so $\delta s = m^\top \delta u$. Then
\[
\delta \mathcal{C}(u,v)
=
\bigl(\delta u-(\delta s)\mathbf{1}_q,\; \delta v+(\delta s)\mathbf{1}_k\bigr).
\]
Pairing this variation with $(\bar u,\bar v)$ and collecting coefficients of $\delta u$ and $\delta v$ yields the stated pullback.
\end{proof}

\begin{lemma}[Gauge Equivariance]
Let $\Gamma_c(u,v)=(u-c\mathbf{1}_q,\; v+c\mathbf{1}_k)$. Then the raw two-half-step map defined by equations~\eqref{eq:u_update_app_unified}--\eqref{eq:v_update_app_unified} is gauge equivariant.
\end{lemma}

\begin{proof}
Adding $c$ to every valid key potential shifts every logsumexp argument in the $u$-update by the same constant, so the update subtracts $c$ from every valid query potential. The $v$-update is identical with rows and columns swapped.
\end{proof}

\begin{proposition}[Centered-Ungauged Forward Equivalence]
Let $(\widehat u^{(t)},\widehat v^{(t)})$ be the ungauged iterates and $(u^{(t)},v^{(t)})$ the iterates produced by the centered implementation, both started from the same centered base pair. Then for every $t\ge 1$ there exists a scalar $c_t$ such that
\[
(u^{(t)},v^{(t)})=\Gamma_{c_t}\bigl(\widehat u^{(t)},\widehat v^{(t)}\bigr).
\]
Consequently the centered and ungauged implementations produce exactly the same staircase transport plans.
\end{proposition}

\begin{proof}
The claim holds at $t=1$ by definition of the centering step. Inductively, if the two previous representatives differ by a gauge transform, then gauge equivariance of the raw refinement step shows that the next raw outputs also differ by a gauge transform; applying $\mathcal{C}$ merely chooses another representative of the same orbit.
\end{proof}

\begin{lemma}[Half-Step Jacobians]
For every refinement step $t\ge 1$,
\begin{align}
    \frac{\partial u_i^{(t)}}{\partial v_j^{(t-1)}} &= -P_{ij}^{(t,t-1)}, \\
    \frac{\partial v_j^{(t)}}{\partial u_i^{(t)}} &= -P_{ij}^{(t,t)}.
\end{align}
\end{lemma}

\begin{proof}
Differentiate the masked logsumexp updates. For example,
\[
\frac{\partial u_i^{(t)}}{\partial v_j^{(t-1)}}
=
-
\frac{\exp(S_{ij}+v_j^{(t-1)})}{\sum_{k:(i,k)\in\Omega}\exp(S_{ik}+v_k^{(t-1)})}
=
-\exp(S_{ij}+u_i^{(t)}+v_j^{(t-1)})
=
-P_{ij}^{(t,t-1)}.
\]
The second identity is identical with rows and columns swapped.
\end{proof}

\paragraph{Proof of Proposition~\ref{prop:tail_adjoint_main}.}
\begin{proof}
The direct terminal transport application contributes $P^{(R,R)}\odot Z$ to $\bar S$ together with
\[
\bar v^{(R)}=g_v,
\qquad
\bar u^{(R)}=g_u-P^{(R,R)}\bar v^{(R)}.
\]
Now fix a reverse step $t\in\{1,\dots,R\}$. Pulling $\bar u^{(t)}$ back through the $u$-half-step uses the Jacobian $-P^{(t,t-1)}$:
\[
\bar v^{(t-1)} = -\bigl(P^{(t,t-1)}\bigr)^\top \bar u^{(t)},
\qquad
\bar S \mathrel{-}= P^{(t,t-1)}\odot\bigl(\bar u^{(t)}\mathbf{1}^\top\bigr).
\]
For $t\ge 2$, pulling $\bar v^{(t-1)}$ back through the $v$-half-step uses the Jacobian $-P^{(t-1,t-1)}$:
\[
\bar u^{(t-1)} = -P^{(t-1,t-1)}\bar v^{(t-1)},
\qquad
\bar S \mathrel{-}= P^{(t-1,t-1)}\odot\bigl(\mathbf{1}\bar v^{(t-1)\top}\bigr).
\]
The recursion stops at $\bar u^{(1)}$ because $u^{(0)}$ belongs to the stopped base solve. Summing the direct terminal term with the reverse-time half-step contributions yields equation~\eqref{eq:generic_score_vjp_unified}.
\end{proof}

\paragraph{Proof of Theorem~\ref{thm:single_base_main}.}
\begin{proof}
For any two staircase plans,
\[
\frac{P_{ij}^{(a,b)}}{P_{ij}^{(2,2)}}
=
\exp(u_i^{(a)}-u_i^{(2)}+v_j^{(b)}-v_j^{(2)}),
\]
which gives the exact reconstruction formula immediately. Substituting these reconstructed plans into the explicit $R=2$ score-space VJP leaves every entry unchanged, so the same backward pass can be evaluated while materializing only $P^{(2,2)}$.
\end{proof}

\section{Proofs for the Dustbin Bridge and Surrogate Bias}
\label{app:bias}

\paragraph{Proof of Theorem~\ref{thm:dustbin_aug}.}
\begin{proof}
After constructing the augmented tensors and masks, the implementation calls the same balanced Sinkhorn forward and backward path as in Section~2.2. Because only the first appended token in each dustbin block is active, every transport entry touching the remaining fillers is identically zero. Because the effective band width is $\widetilde W=\max(W,L)$, every active base token can reach the active dustbin token on the opposite side. Thus the implemented path is exactly the balanced fixed-support surrogate on an enlarged state space.
\end{proof}

\paragraph{Proof of Corollary~\ref{cor:dustbin_r2}.}
\begin{proof}
Apply Proposition~\ref{prop:tail_adjoint_main} and Theorem~\ref{thm:single_base_main} to the augmented tensors and support from Theorem~\ref{thm:dustbin_aug}. No new algebra is introduced after the augmentation itself.
\end{proof}

\paragraph{Proof of Theorem~\ref{thm:surrogate_bias}.}
\begin{proof}
By the chain rule,
\[
\nabla_x\Psi_x^R(z_T(x))
=
\partial_x \Psi_x^R(z)\big|_{z=z_T(x)}
+
D_z\Psi_x^R(z_T(x))\,D_x z_T(x),
\]
which gives the exact bias identity after subtracting the partial derivative term.

Next,
\[
D_z\Psi_x^R(z_T(x))
=
D_z A_x(z_{T+R}(x))
\prod_{r=0}^{R-1} D_z\Phi_x(z_{T+r}(x)),
\]
so
\[
\|D_z\Psi_x^R(z_T(x))\|
\le
L_A \rho^R.
\]
Differentiate the base recursion $z_{t+1}(x)=\Phi_x(z_t(x))$:
\[
D_x z_{t+1}(x)
=
D_x\Phi_x(z_t(x))
+
D_z\Phi_x(z_t(x))\,D_x z_t(x).
\]
Because $z_0$ is independent of $x$, repeated application of the norm bounds yields
\[
\|D_x z_T(x)\|
\le
L_\Phi \sum_{j=0}^{T-1}\rho^j
\le
\frac{L_\Phi}{1-\rho}.
\]
Combining the two estimates with the exact bias identity gives
\[
\bigl\|g_{\mathrm{full}}^{T,R}(x)-g_{\mathrm{surr}}^{T,R}(x)\bigr\|
\le
\frac{L_A L_\Phi}{1-\rho}\,\rho^R.
\qedhere
\]
\end{proof}

\section{Experimental Protocol Notes}

The current manuscript uses three empirical layers:
\begin{enumerate}
    \item exactness checks against exact autodiff on small masked problems;
    \item a four-configuration TPU v6e-8 Pfam screen across balanced and dustbin-augmented paths;
    \item a promoted long-budget balanced run for the theorem-backed $R=2$ path.
\end{enumerate}

The present cloud-aligned gap-aware path is the dustbin-augmented balanced surrogate. Genuinely KL-unbalanced transport and differentiable support geometry remain natural follow-on directions, but they are not needed for the main claim carried by this manuscript.

\end{document}
